%==============================================================
%  schemes.tex
%  hasklib — mathematical companion / derivations
%
%  Numerical schemes for simulating an Ito SDE: Euler-Maruyama now;
%  Milstein to be added once Milstein.hpp/.cpp exist (same file,
%  code and math built together, matching library convention).
%
%  Cites ito_foundations.tex for the SDE framework and the
%  Gaussian-increment property of Brownian motion.
%
%  Scope discipline: only Euler-Maruyama is derived here, because
%  only Euler-Maruyama exists as verified code so far
%  (EulerMaruyama::step / simulate_terminal, tested against GBM's
%  exact sampler in test_stochastic.cpp).
%==============================================================
\documentclass[11pt]{article}

\usepackage[utf8]{inputenc}
\usepackage[T1]{fontenc}
\usepackage{amsmath, amssymb, amsthm}
\usepackage{mathtools}
\usepackage[a4paper, margin=1in]{geometry}
\usepackage{hyperref}
\usepackage{cleveref}
\usepackage{bm}
\usepackage{xcolor}
\usepackage{harvard}   % Harvard author-date referencing (with agsm.bst)

%---------------- notation macros (matching ito_foundations.tex) ----------------
\newcommand{\Prob}{\mathbb{P}}
\newcommand{\Qmeas}{\mathbb{Q}}
\newcommand{\E}{\mathbb{E}}
\newcommand{\Var}{\operatorname{Var}}
\newcommand{\Cov}{\operatorname{Cov}}
\newcommand{\filt}{\mathcal{F}}
\newcommand{\Om}{\Omega}
\newcommand{\dd}{\mathrm{d}}
\newcommand{\dW}{\dd W_t}
\newcommand{\dt}{\dd t}
\newcommand{\R}{\mathbb{R}}
\newcommand{\drift}{a}
\newcommand{\diff}{b}

\theoremstyle{plain}
\newtheorem{theorem}{Theorem}
\newtheorem{lemma}{Lemma}
\newtheorem{proposition}{Proposition}

\theoremstyle{definition}
\newtheorem{definition}{Definition}
\newtheorem{assumption}{Assumption}

\theoremstyle{remark}
\newtheorem{remark}{Remark}

\newcommand{\TODO}[1]{\medskip\noindent\textbf{[\textcolor{red}{TO DERIVE}: #1]}\medskip}

%==============================================================
\title{Numerical Schemes for SDE Simulation: Euler--Maruyama\\
       \large \texttt{hasklib} mathematical companion}
\author{Hubert}
\date{\today}
%==============================================================

\begin{document}
\maketitle

\begin{abstract}
This note derives the Euler--Maruyama scheme for simulating an
It\^o SDE, connects it directly to
\texttt{EulerMaruyama::step}/\texttt{simulate\_terminal}, and states
its strong-order convergence result. It builds on the SDE framework
and Brownian motion properties established in
\texttt{ito\_foundations.tex}. The Milstein scheme is deliberately
omitted here and will be appended once \texttt{Milstein.hpp/.cpp}
exist, keeping code and derivation paired throughout the library.
\end{abstract}

%--------------------------------------------------------------
\section{Setup: the problem discretisation solves}
%--------------------------------------------------------------

Recall from \texttt{ito\_foundations.tex} the general It\^o SDE
\[
   \dd X_t = \drift(t,X_t)\,\dt + \diff(t,X_t)\,\dW.
\]
Integrating exactly over one interval $[t,\,t+\Delta t]$ gives
\begin{equation}\label{eq:exact-increment}
   X_{t+\Delta t} - X_t
   = \int_t^{t+\Delta t} \drift(s,X_s)\,\dd s
   + \int_t^{t+\Delta t} \diff(s,X_s)\,\dd W_s.
\end{equation}
This is exact, but neither integral can be evaluated in closed form
in general, since $\drift$ and $\diff$ depend on $X_s$, which itself
moves randomly across the interval. Simulation requires an
approximation of \eqref{eq:exact-increment}.

%--------------------------------------------------------------
\section{The Euler--Maruyama scheme}
%--------------------------------------------------------------

\begin{definition}[Euler--Maruyama update]\label{def:em}
The Euler--Maruyama approximation freezes $\drift$ and $\diff$ at
their values at the start of the interval and approximates
\eqref{eq:exact-increment} by
\begin{equation}\label{eq:em-approx}
   \int_t^{t+\Delta t}\drift(s,X_s)\,\dd s \approx \drift(t,X_t)\,\Delta t,
   \qquad
   \int_t^{t+\Delta t}\diff(s,X_s)\,\dd W_s \approx \diff(t,X_t)\big(W_{t+\Delta t}-W_t\big).
\end{equation}
\end{definition}

\begin{remark}[Why freezing at $t$: the defining choice of an explicit scheme]
By the mean value theorem for integrals, there exists some
$\xi\in[t,t+\Delta t]$ (depending on the whole path) with
\[
   \int_t^{t+\Delta t}\drift(s,X_s)\,\dd s = \drift(\xi,X_\xi)\,\Delta t
   \quad\text{exactly.}
\]
This is an exact statement but a useless one, since $\xi$ and $X_\xi$
are unknown until the interval has already been simulated. Euler--
Maruyama's approximation is to replace the unknown midpoint value
$\drift(\xi,X_\xi)$ with the \emph{known} left-endpoint value
$\drift(t,X_t)$ --- valid to leading order because $\drift$ is
continuous and the interval is short, so $\drift(\xi,X_\xi)$ differs
from $\drift(t,X_t)$ by an amount that shrinks as $\Delta t\to0$
(quantified in \Cref{sec:local-error} below). The diffusion
approximation in \eqref{eq:em-approx} is the same idea applied to the
stochastic integral: if $\diff(s,X_s)$ were genuinely constant over
$[t,t+\Delta t]$, the stochastic integral would equal
$\diff(t,X_t)(W_{t+\Delta t}-W_t)$ \emph{exactly}, so freezing it at
the left endpoint is again a leading-order replacement of an unknown
quantity by a known one.

The choice of \emph{which} endpoint to freeze at is what separates
explicit from implicit schemes. Freezing at $t$ means every quantity
on the right of \eqref{eq:em-approx} is already known once we reach
time $t$, so $X_{t+\Delta t}$ is computed by a single formula with no
equation to solve. Freezing instead at the right endpoint
$t+\Delta t$ (backward Euler) or at the midpoint would put the unknown
$X_{t+\Delta t}$ (or an average involving it) on both sides of the
update, requiring a root-find at every step. Explicitness is a
computational convenience, not a mathematical requirement of the
discretisation idea itself --- it is a choice about which known
quantity approximates $\drift(\xi,X_\xi)$, and hasklib's Euler--
Maruyama always chooses the one requiring no further work.
\end{remark}

Using the Gaussian-increment property of Brownian motion
(\texttt{ito\_foundations.tex}, Definition~1), write
$W_{t+\Delta t}-W_t = \sqrt{\Delta t}\,Z$ with $Z\sim\mathcal N(0,1)$.
Substituting into \eqref{eq:em-approx} and \eqref{eq:exact-increment}
gives the update rule:
\begin{equation}\label{eq:em-update}
   X_{t+\Delta t} \;\approx\; X_t + \drift(t,X_t)\,\Delta t
                    + \diff(t,X_t)\sqrt{\Delta t}\,Z.
\end{equation}

\begin{remark}[Correspondence with the code]\label{rem:em-code}
Equation~\eqref{eq:em-update} is exactly \texttt{EulerMaruyama::step}:
\begin{verbatim}
double EulerMaruyama::step(const StochasticProcess& process, double t,
                            double x, double dt, double z) const
{
    double a = process.drift(t, x);
    double b = process.diffusion(t, x);
    return x + a * dt + b * std::sqrt(dt) * z;
}
\end{verbatim}
Chaining \eqref{eq:em-update} across $N$ equal sub-intervals of
$[0,T]$, with a fresh independent $Z_i$ drawn at each step (the
independent-increments property of Brownian motion), is exactly
\texttt{EulerMaruyama::simulate\_terminal}.
\end{remark}

%--------------------------------------------------------------
\section{Why the approximation improves as $\Delta t\to 0$}\label{sec:local-error}
%--------------------------------------------------------------

The error introduced by \Cref{def:em} comes entirely from treating
$\drift(s,X_s)$ and $\diff(s,X_s)$ as constant over $[t,t+\Delta t]$,
when in fact $X_s$ moves within the interval. By the Brownian scaling
argument in \texttt{ito\_foundations.tex} (the nowhere-differentiability
discussion), that movement is of order $\sqrt{\Delta t}$, so the
error made by freezing the coefficients should itself shrink as
$\Delta t \to 0$.

\begin{proposition}[One-step (local) mean-square error]\label{prop:local-error}
Assume $\drift,\diff$ are Lipschitz and continuously differentiable in
$x$. Writing $\varepsilon$ for the difference between the true
increment \eqref{eq:exact-increment} and the Euler--Maruyama
approximation \eqref{eq:em-approx} over one step of length
$\Delta t$, started from the same $X_t$,
\[
   \E\big[\varepsilon^2\big] = O\big((\Delta t)^2\big),
   \qquad\text{i.e. the root-mean-square local error is } O(\Delta t).
\]
\end{proposition}

\begin{proof}[Heuristic derivation]
Split $\varepsilon$ into its drift and diffusion contributions,
$\varepsilon = \varepsilon_{\drift} + \varepsilon_{\diff}$, where
\[
   \varepsilon_{\drift} = \int_t^{t+\Delta t}\!\big(\drift(s,X_s)-\drift(t,X_t)\big)\dd s,
   \qquad
   \varepsilon_{\diff} = \int_t^{t+\Delta t}\!\big(\diff(s,X_s)-\diff(t,X_t)\big)\dd W_s.
\]
For $s\in[t,t+\Delta t]$, the true solution satisfies
$X_s - X_t = \drift(t,X_t)(s-t) + \diff(t,X_t)(W_s-W_t) + \cdots$; the
first term is $O(\Delta t)$ and the second is $O(\sqrt{\Delta t})$ in
an $L^2$ sense (its variance is $\diff(t,X_t)^2(s-t)$), so the
Brownian term dominates and $X_s-X_t = O(\sqrt{\Delta t})$. By the
Lipschitz/differentiability assumption, both coefficient differences
inherit this order:
\[
   \drift(s,X_s)-\drift(t,X_t) = O(\sqrt{\Delta t}), \qquad
   \diff(s,X_s)-\diff(t,X_t)   = O(\sqrt{\Delta t}).
\]
\emph{Drift term.} Integrating a quantity of order $\sqrt{s-t}$ over
$s\in[t,t+\Delta t]$ gives $\int_0^{\Delta t}\sqrt{u}\,\dd u =
\tfrac23(\Delta t)^{3/2}$, so $\varepsilon_{\drift}=O\big((\Delta
t)^{3/2}\big)$ pathwise --- already smaller than the diffusion
contribution below, so it does not set the overall order.

\emph{Diffusion term.} Here the integral is stochastic, so its size is
governed by the (adapted) It\^o isometry --- the same fact used for
the Moments of a Wiener Integral lemma in \texttt{ou.tex}, but for an
integrand that is random and adapted rather than deterministic; see
\citeasnoun{shreve2004} or \citeasnoun{oksendal2003} for this more
general form:
\[
   \E\big[\varepsilon_{\diff}^2\big]
     = \E\left[\int_t^{t+\Delta t}\!\big(\diff(s,X_s)-\diff(t,X_t)\big)^2\dd s\right]
     \approx \int_t^{t+\Delta t} O(s-t)\,\dd s = O\big((\Delta t)^2\big),
\]
using $\E\big[(\diff(s,X_s)-\diff(t,X_t))^2\big]=O(s-t)$ from the
$L^2$ bound above. So $\varepsilon_{\diff}$ has mean-square
$O((\Delta t)^2)$, i.e.\ root-mean-square $O(\Delta t)$, which
dominates the drift contribution and gives the proposition.
\end{proof}

\begin{remark}
This is the mechanism, not yet the global result: a per-step error
of root-mean-square order $\Delta t$ does not immediately give a
global error of the same order, because the errors from different
steps compound through the process's dependence on $X_t$ rather than
simply adding. That compounding is the content of
\Cref{thm:em-order} below.
\end{remark}

%--------------------------------------------------------------
\section{Strong convergence}
%--------------------------------------------------------------

\begin{definition}[Strong convergence order]
A scheme has \emph{strong order} $\gamma$ if there exists $C>0$,
independent of the step size $\Delta t$, such that
\[
   \E\Big[\,\big|\hat X_T - X_T\big|\,\Big] \;\le\; C\,(\Delta t)^{\gamma},
\]
where $\hat X_T$ is the scheme's approximation and $X_T$ is the true
solution, evaluated pathwise on the same Brownian driving noise.
\end{definition}

\begin{theorem}[Strong order of Euler--Maruyama]\label{thm:em-order}
Under standard regularity conditions on $\drift$ and $\diff$
(Lipschitz continuity and linear growth, as in
\texttt{ito\_foundations.tex}, Assumption~1), the Euler--Maruyama
scheme has strong order $\gamma = \tfrac12$.
\end{theorem}

\begin{proof}[Proof strategy]
\Cref{prop:local-error} gives a per-step root-mean-square error of
order $\Delta t$. Simulating to time $T$ takes $N=T/\Delta t$ such
steps, and the natural (but wrong) guess is that independent
per-step errors of size $\Delta t$ accumulate to a global error of
size $N\cdot\Delta t = T$ --- a constant, not something that shrinks.
That guess is wrong because it ignores that each step's error is fed
back through $\drift$ and $\diff$ at the start of the \emph{next}
step: an error made at step $i$ perturbs $X_{t_i}$, which perturbs
every subsequent coefficient evaluation. Controlling this requires a
discrete Gr\"onwall-type argument: the Lipschitz bound on $\drift$
and $\diff$ shows that an error present at step $i$ can grow by at
most a bounded multiplicative factor per subsequent step, so the
$N$ accumulated local errors combine as a sum dominated by the
largest few terms rather than by their count. Carrying this through
converts the local root-mean-square bound $O(\Delta t)$ into a global
strong error of $O\big(\sqrt{N}\cdot\Delta t\big) = O\big(\sqrt{T/\Delta
t}\cdot\Delta t\big) = O(\sqrt{\Delta t})$, i.e.\ strong order
$\tfrac12$. The technical machinery (the precise Gr\"onwall estimate
and the regularity conditions under which it holds) is in
\citeasnoun{kloeden1992}, Ch.~9--10; we cite the result rather than
reproduce the estimate.
\end{proof}

\begin{remark}[Correspondence with the code]
This is precisely what \texttt{test\_stochastic.cpp}'s convergence
block checks empirically: running \texttt{simulate\_terminal} on
\texttt{GBM} at increasing step counts $N \in \{10,50,200,1000\}$
and observing the empirical terminal moments approach the exact
moments from \texttt{GBM::sample\_terminal} as $\Delta t = T/N \to 0$.
\Cref{thm:em-order} is the theoretical statement of which the test's
shrinking printed errors are a numerical demonstration.
\end{remark}

\begin{remark}[Contrast with the exact sampler]
For GBM specifically, \texttt{GBM::sample\_terminal} has \emph{zero}
discretisation error, because the SDE admits a closed-form solution
via It\^o's lemma (see \texttt{gbm.tex}). Euler--Maruyama is an
approximation of the \emph{same} SDE and is only needed in general
because most processes --- and, in particular, any path-dependent
payoff even for GBM --- do not admit such a closed form.
\end{remark}

%--------------------------------------------------------------
\section*{Milstein scheme --- reserved}
%--------------------------------------------------------------
\TODO{add the Milstein derivation, strong order 1.0 result, and its
correspondence remark once \texttt{Milstein.hpp/.cpp} are written and
tested. Note the diffusion-derivative term $b\,b'\,(Z^2-1)/2$ vanishes
exactly when $b$ is constant (ABM, OU), giving a clean regression
check that Milstein degenerates to Euler on those processes.}

%--------------------------------------------------------------
\bibliographystyle{agsm}
\bibliography{refs}

\end{document}